%  Preprints.org template — MDPI mdpi.cls
\documentclass[preprints,article,accept,pdftex,oneauthor]{Definitions/mdpi}

% MDPI internal commands — do not modify
\firstpage{1}
\makeatletter
\setcounter{page}{\@firstpage}
\makeatother
\pubvolume{1}
\issuenum{1}
\articlenumber{0}
\pubyear{2026}
\copyrightyear{2026}
\datereceived{ }
\daterevised{ }
\dateaccepted{ }
\datepublished{ }
\hreflink{https://doi.org/}

%=================================================================
\Title{Foundational Reasoning and Feedback Protocol (FRFP): An Updated Formal Specification}
\TitleCitation{Foundational Reasoning and Feedback Protocol (FRFP): An Updated Formal Specification}

\Author{Vijaya Nadendla $^{1,\dagger}$}
\AuthorNames{Vijaya Nadendla}
\isAPAStyle{%
  \AuthorCitation{Nadendla, V.}
}{%
  \isChicagoStyle{%
    \AuthorCitation{Nadendla, Vijaya.}
  }{%
    \AuthorCitation{Nadendla, V.}
  }
}

\address{$^{1}$ \quad Independent Researcher}
\firstnote{Current address: 51 Linwood Dr, Monroe Township, NJ 08831.}

\abstract{The Foundational Reasoning and Feedback Protocol (FRFP) is a formal framework for
structuring reasoning and evaluation in long-horizon human--AI workflows. It distinguishes between
explicit processes, which operate on representations and are mechanically reproducible, and tacit
processes, which determine correctness, validity, and acceptance. The framework enforces a one-way
boundary between these domains, under which explicit computation produces artifacts that must be
evaluated through tacit judgment.
This paper presents the FRFP theory in the form established by complete machine-checked
formalization in Lean~4. The formalization yields four refinements to the published theory:
(1)~Basis reduction: the original 11 premises reduce to 7 independent axioms, with four
(ETS, RB, AE, MD) now derivable theorems; (2)~Precision corrections: six definitions required
strengthening, notably explicit credence bounds, clamped degradation, and normal-form uniqueness
modulo navigation equivalence; (3)~New formal structures: three types implicit in the original
paper are made explicit --- the runToOmega bijection, the CorrectnessJudgment type, and
lifted-action commutativity; (4)~Novel theorems: three non-trivial results emerge from proof
obligations --- a canonical first-unresolved-witness theorem, a governance-safe trace-extension
theorem, and a protocol normal-form decomposition theorem.
Where the published paper and Lean code diverge, the Lean version is stated explicitly and
documented.}

\keyword{human--AI collaboration; epistemic governance; tacit knowledge; explicit--tacit separation;
AI safety; long-horizon reasoning; category theory; sociotechnical systems}

%=================================================================
% Additional packages (not already loaded by mdpi.cls)
\usepackage{mathtools}
\usepackage{stmaryrd}
\usepackage{longtable}
\usepackage{listings}
\usepackage{tcolorbox}
\usepackage{mdframed}
\usepackage{adjustbox}

\setlength{\emergencystretch}{3em}
\AtBeginEnvironment{tabular}{\small}
\BeforeBeginEnvironment{tabular}{\begin{adjustbox}{max width=\textwidth,center}}
\AfterEndEnvironment{tabular}{\end{adjustbox}}

% --- Theorem environments ---
% mdpi.cls already defines: Theorem, Lemma, Corollary, Proposition,
%   Definition, Example, Remark, Notation, Assumption (capitalized).
% Provide lowercase aliases for compatibility with chapter files.
\newenvironment{theorem}{\begin{Theorem}}{\end{Theorem}}
\newenvironment{lemma}{\begin{Lemma}}{\end{Lemma}}
\newenvironment{corollary}{\begin{Corollary}}{\end{Corollary}}
\newenvironment{proposition}{\begin{Proposition}}{\end{Proposition}}
\newenvironment{definition}{\begin{Definition}}{\end{Definition}}
\newenvironment{remark}{\begin{Remark}}{\end{Remark}}
\newenvironment{example}{\begin{Example}}{\end{Example}}
% axiom_env and noteenv not in mdpi.cls — define separately
\newcounter{axiomctr}
\newtheorem{axiom_env}{Axiom}[section]
\newcounter{notectr}
\newtheorem{noteenv}{Note}[section]
% Alias \begin{note} to noteenv
\let\note\noteenv
\let\endnote\endnoteenv

% --- Lean reference ---
\newcommand{\leanref}[1]{\textbf{[Lean:} \nolinkurl{#1}\textbf{]}}
\newcommand{\leanbox}[1]{%
  \begin{mdframed}[backgroundcolor=gray!10, linewidth=0.5pt]
    \small\textbf{Lean:} \texttt{\detokenize{#1}}
  \end{mdframed}}

% --- Math shorthands ---
\newcommand{\Czero}{\mathbf{C}_0}
\newcommand{\Ecat}{\mathbf{E}_{\mathrm{cat}}}
\newcommand{\Tcat}{\mathbf{T}_{\mathrm{cat}}}
\newcommand{\Eobj}{E_0}
\newcommand{\Tobj}{T_0}
\newcommand{\runToOmega}{\mathrm{runTo}\omega}
\newcommand{\omegaToRun}{\omega\mathrm{ToRun}}
\newcommand{\navequiv}{\equiv_\nabla}
\newcommand{\RI}{\mathrm{RI}}
\newcommand{\EC}{\mathrm{EC}}
\newcommand{\ED}{\mathrm{ED}}
\newcommand{\RB}{\mathrm{RB}}
\newcommand{\TE}{\mathrm{TE}}
\newcommand{\HFD}{\mathrm{HFD}}
\newcommand{\obs}{\mathrm{obs}}
\newcommand{\Sem}{\mathrm{Sem}}

% --- Lean listings style ---
\lstdefinelanguage{Lean4}{
  keywords={def,theorem,lemma,axiom,structure,inductive,where,match,with,
            fun,let,in,have,show,by,intro,apply,exact,simp,rfl,cases,induction,
            forall,exists,if,then,else,namespace,end,import,open,section,
            class,instance,variable,noncomputable,opaque},
  keywordstyle=\color{blue!70!black}\bfseries,
  comment=[l]{--},
  morecomment=[s]{/-}{-/},
  commentstyle=\color{green!50!black}\itshape,
  stringstyle=\color{red!70!black},
  sensitive=true,
  basicstyle=\ttfamily\small,
  breaklines=true,
  frame=single,
  frameround=tttt,
  backgroundcolor=\color{gray!5},
  columns=fullflexible,
  keepspaces=true,
  upquote=true,
  literate=
    {<=}{{$\le$}}2
    {>=}{{$\ge$}}2
    {->}{{$\to$}}2
    {≤}{{$\le$}}1
    {≥}{{$\ge$}}1
    {→}{{$\to$}}1
    {↔}{{$\leftrightarrow$}}1
    {∀}{{$\forall$}}1
    {∃}{{$\exists$}}1
    {∧}{{$\land$}}1
    {∨}{{$\lor$}}1
    {¬}{{$\neg$}}1
    {§}{{\S}}1
    {·}{{$\cdot$}}1
    {ω}{{$\omega$}}1
    {Ω}{{$\Omega$}}1
    {⟨}{{$\langle$}}1
    {⟩}{{$\rangle$}}1,
}
\lstset{language=Lean4,breaklines=true,breakatwhitespace=false,%
  postbreak=\mbox{\textcolor{gray}{$\hookrightarrow$}\space}}

% --- Correction box ---
\newenvironment{correction}[1][Lean Correction]{%
  \begin{tcolorbox}[colback=orange!10,colframe=orange!50!black,title={\textbf{#1}}]
}{%
  \end{tcolorbox}
}

%=================================================================
\begin{document}
%=================================================================

\maketitle

\newpage
\tableofcontents
\newpage

% ===================================================================
\section*{Notation}
% ===================================================================

\begin{center}
\begin{tabular}{ll}
\toprule
Symbol & Meaning \\
\midrule
$\Czero$ & The kernel category with objects $\{\varnothing, \Eobj, \Tobj\}$ \\
$\Eobj$ & The explicit object \\
$\Tobj$ & The tacit object \\
$\Ecat$ & The explicit subcategory (morphisms targeting $\Eobj$) \\
$\Tcat$ & The tacit subcategory (morphisms targeting $\Tobj$) \\
$\Pi$ & A pipeline (morphism in $\Czero$) \\
$\mathcal{P}$ & A protocol trace (list of step records) \\
$\sigma$ & A configuration (object in the Grothendieck construction) \\
$\navequiv$ & Navigation equivalence relation \\
$\delta$ & Degradation operator on credence \\
$S(n)$ & Survival function at step $n$ \\
$h(n)$ & Hazard rate at step $n$ \\
$\omega$ & A stochastic trajectory \\
$\Omega$ & Sample space of stochastic trajectories \\
\bottomrule
\end{tabular}
\end{center}

Throughout the paper, Lean theorems and definitions are referenced using the pattern
\leanref{module.theorem\_name}. All theorems marked \leanref{\ldots} have zero
\texttt{sorry} and are accepted by the Lean~4.29.0 kernel.

\newpage

% ===================================================================
% Chapters
% ===================================================================
\section{Introduction}
\label{ch:intro}

The Foundational Reasoning and Feedback Protocol (FRFP) is a formal framework for structuring
reasoning and evaluation in systems that combine machine computation with human judgment. It
is based on a fundamental distinction between \emph{explicit} processes, which operate on
representations and can be executed mechanically, and \emph{tacit} processes, which determine
correctness, validity, and acceptance.

Earlier presentations of FRFP introduced this structure as a conceptual and theoretical framework
for long-horizon human--AI collaboration. This paper provides an updated and consolidated formal
specification of the framework. The presentation reflects refinements arising from complete
formalization, including simplification of the primitive basis, clarification of key definitions,
and strengthening of structural constraints.

The purpose of this document is to present the current form of FRFP as a self-contained system.
The focus is on precise definitions and derivable properties rather than on formal verification or proof
mechanics. Where earlier versions differ, the refined formulation given here is taken as authoritative.

The remainder of the paper develops the framework starting from a minimal categorical kernel
that captures the separation between explicit and tacit domains.

\subsection*{Validation Boundary: Lean and Human Roles}

This theory paper follows the same validation boundary used in the Lean companion.
Lean is the correctness oracle for formal derivations: once a theorem is accepted by
Lean, no human re-check of proof steps is required. Human validation serves a
\emph{different} role: it governs commitments.

Concretely:
\begin{itemize}
  \item \textbf{Lean verifies proofs}: derivations are correct relative to stated definitions and axioms.
  \item \textbf{Humans validate commitments}: statement intent, premise legitimacy, and publication scope.
\end{itemize}

Under FRFP, this distinction is structural rather than optional: explicit derivation is
machine-checkable, while normative acceptance remains a tacit-space responsibility.

\subsection*{The FRFP Document Ecosystem}

This paper is part of a broader research program organized into three layers. Understanding
where this document fits will help the reader navigate the full body of work.

\textbf{Theory layer.} The mathematical foundations of FRFP are developed across three documents.
\textit{FRFPPaper Vol.~1} introduced the explicit--tacit separation, the representational boundary,
and the formal protocol constraints. The present document supersedes that account with a
consolidated and corrected specification, incorporating refinements that emerged from complete
formalization. For readers who require a pedagogical introduction to the underlying mathematics,
the \textit{FRFP Companion Tutorial} provides accessible explanations of the key structures
before any formal development. The \textit{Lean Formalization} companion paper records the
machine-checked verification of this theory in Lean~4, including 269 theorems, a basis reduction
from 11 premises to 7, and novel results not present in earlier versions.

\textbf{Engineering layer.} \textit{FRFPPaper Vol.~2} applies the theoretical framework to a
corpus of canonical socio-technical case studies. The \textit{FRFP Technical Specification}
translates the abstract protocol constraints into implementable procedures for LLM-based
reasoning pipelines. A series of \textit{FRFP White Papers} (e.g., \textit{Beyond Prompt
Engineering}) provides engineering-facing interpretations of the protocol and its practical
consequences.

\textbf{Evaluation layer.} \textit{FRFP Vol.~3: Experimental Methods for Evaluating AI Agents} defines a
general controlled-evaluation framework for AI reasoning systems. The \textit{FRFP vs.\ Base
Survey} applies that framework to empirically assess protocol-guided reasoning against baseline
LLM outputs, using normalized one-page analytical summaries as the evaluation artifacts for
blinded human raters.

\medskip
\noindent\textbf{Recommended reading path.} A reader new to FRFP should begin with
\textit{FRFPPaper Vol.~1}, which provides the motivation, conceptual development, and
theoretical arguments underlying the framework. The present document should then be read
as the authoritative formal specification: it refines and consolidates the definitions from
Vol.~1 but does not repeat its motivating discussion. Readers who find the mathematics
unfamiliar should consult the FRFP Companion Tutorial before or alongside this document.
After this document, readers interested in machine-checked proofs should proceed to the
Lean Formalization companion. Readers whose primary interest is deployment or empirical
evaluation should continue to the Engineering and Evaluation layers in sequence.

Practitioners and engineers whose primary interest is implementation rather than formal
theory may proceed directly to \textit{FRFPPaper Vol.~2}, which provides an
engineering-oriented exposition of the framework through detailed socio-technical case
studies, and to the \textit{FRFP Technical Specification} for operational protocol rules.

\newpage
% ===================================================================
\section{The Kernel Category}
\label{ch:kernel}
% ===================================================================

\subsection{Overview}

FRFP is built on a minimal three-object category $\Czero$ whose structure encodes
the fundamental distinction between explicit (machine-manipulable) and tacit
(human correctness-bearing) computation. This chapter defines $\Czero$, its six
primitive morphisms, their typing rules, and the core theorems that follow directly
from the definitions --- without requiring any additional axioms.

% -------------------------------------------------------------------
\subsection{Objects}
% -------------------------------------------------------------------

\begin{definition}[Kernel category objects]
The kernel category $\Czero$ has exactly three objects:
\[
  \mathrm{Ob}(\Czero) = \{\varnothing,\ \Eobj,\ \Tobj\}
\]
\end{definition}

\begin{center}
\begin{tabular}{lll}
\toprule
Object & Name & Interpretation \\
\midrule
$\varnothing$ & Empty & The absent or uninitialized explicit state \\
$\Eobj$ & Explicit object & Machine-manipulable representations \\
$\Tobj$ & Tacit object & Human correctness-bearing knowledge \\
\bottomrule
\end{tabular}
\end{center}

The object set is closed. No extensions to $\Czero$ are permitted.
\leanref{Frfp.Core.Kernel.Object}

The objects are sometimes written with subscript notation: $\Eobj$ (the explicit
\emph{object}) is distinct from $\Ecat$ (the \emph{category} of explicit
morphisms). This distinction is enforced in the Lean formalization by separate
types.

% -------------------------------------------------------------------
\subsection{Primitive Morphisms}
% -------------------------------------------------------------------

\begin{definition}[Primitives]
The six primitive morphisms of $\Czero$ are:
\end{definition}

\begin{center}
\begin{tabular}{llll>{\raggedright\arraybackslash}p{4.5cm}}
\toprule
Primitive & Source & Target & Name & Role \\
\midrule
$\RI$ & $\varnothing$ & $\Eobj$ & Representation Initiation & Creates an explicit representation \\
$\EC$ & $\Eobj$ & $\Eobj$ & Explicit Computation & AI processing within explicit space \\
$\ED$ & $\Eobj$ & $\Eobj$ & Explicit Diagnostics & AI checking/verification within explicit space \\
$\RB$ & $\Eobj$ & $\Tobj$ & Representational Backflow & The unique boundary crossing morphism \\
$\TE$ & $\Tobj$ & $\Tobj$ & Tacit Evaluation & Human evaluation within tacit space \\
$\HFD$ & $\Tobj$ & $\Tobj$ & Human Final Decision & Human closure within tacit space \\
\bottomrule
\end{tabular}
\end{center}

\leanref{Frfp.Core.Kernel.Primitive}

The primitive set is a closed enumeration. No new primitives may be added without
a complete re-audit of all downstream theorems.

% -------------------------------------------------------------------
\subsection{Subcategories}
% -------------------------------------------------------------------

\begin{definition}[Explicit subcategory]
$\Ecat$ is the full subcategory of morphisms targeting $\Eobj$:
\[
  \Ecat = \{p : \mathrm{Primitive} \mid p.\mathrm{target} = \Eobj\}
        = \{\RI, \EC, \ED\}
\]
\end{definition}

\begin{definition}[Tacit subcategory]
$\Tcat$ is the full subcategory of morphisms targeting $\Tobj$:
\[
  \Tcat = \{p : \mathrm{Primitive} \mid p.\mathrm{target} = \Tobj\}
        = \{\RB, \TE, \HFD\}
\]
\end{definition}

\leanref{Frfp.Core.Kernel.Ecat}, \leanref{Frfp.Core.Kernel.Tcat}

The partition $\Ecat \cup \Tcat$ covers all six primitives and is disjoint.

% -------------------------------------------------------------------
\subsection{Core Typing Theorems}
% -------------------------------------------------------------------

The following theorems are proved directly from the primitive definitions --- no
additional axioms are required.

\begin{theorem}[Explicit--Tacit Separation]
There is no primitive morphism from $\Tobj$ to $\Eobj$:
\[
  \forall p : \mathrm{Primitive},\
  \neg(p.\mathrm{source} = \Tobj \wedge p.\mathrm{target} = \Eobj)
\]
\end{theorem}
\begin{proof}
By case analysis on the six primitives. None has source $\Tobj$ and target
$\Eobj$. \qed
\end{proof}
\leanref{Frfp.Core.Kernel.no\_morphism\_T0\_to\_E0}

\begin{remark}
In the original published theory, ETS was listed as a primitive axiom. The Lean
formalization shows it is a consequence of the primitive typing. See
Chapter~\ref{ch:basis}.
\end{remark}

\begin{theorem}[Uniqueness of the boundary morphism]
$\RB$ is the unique primitive crossing from $\Eobj$ to $\Tobj$:
\[
  \forall p : \mathrm{Primitive},\
  (p.\mathrm{source} = \Eobj \wedge p.\mathrm{target} = \Tobj) \Rightarrow p = \RB
\]
\end{theorem}
\begin{proof}
By case analysis on the six primitives. Only $\RB$ has source $\Eobj$ and target
$\Tobj$. \qed
\end{proof}
\leanref{Frfp.Core.Kernel.RB\_unique\_boundary}

\begin{theorem}[AI executability partition]
Define \texttt{is\_AI\_executable} as $\{\RI, \EC, \ED\}$. Then:
\[
  \forall p,\
  \textit{is\_AI\_executable}(p)
  \Rightarrow (p.\mathrm{source} \in \{\varnothing, \Eobj\})
              \wedge p.\mathrm{target} = \Eobj
\]
\end{theorem}
\leanref{Frfp.Core.Kernel.AI\_preserves\_explicit}

\begin{corollary}[Human-exclusive tacit access]
$\RB$, $\TE$, and $\HFD$ are never AI-executable and always operate in or
enter tacit space.
\leanref{Frfp.Core.Kernel.is\_human\_only}
\end{corollary}

% -------------------------------------------------------------------
\subsection{Pipelines}
% -------------------------------------------------------------------

\begin{definition}[Pipeline]
A pipeline $\Pi$ is a morphism in $\Czero$: a triple $(\mathrm{src},
\mathrm{tgt}, p)$ where $p$ is the underlying primitive.
\end{definition}

\begin{center}
\begin{tabular}{llll}
\toprule
Pipeline type & Source & Target & Underlying primitives \\
\midrule
Explicit pipeline & $\varnothing$ or $\Eobj$ & $\Eobj$ & $\RI$, $\EC$, $\ED$ \\
Boundary pipeline & $\Eobj$ & $\Tobj$ & $\RB$ (unique) \\
Tacit pipeline & $\Tobj$ & $\Tobj$ & $\TE$, $\HFD$ \\
\bottomrule
\end{tabular}
\end{center}

\begin{definition}[Pipeline composition]
For pipelines $\Pi_1 : X \to Y$ and $\Pi_2 : Y \to Z$, their composition
$\Pi_2 \circ \Pi_1 : X \to Z$ is defined when
$\Pi_1.\mathrm{target} = \Pi_2.\mathrm{source}$.
\leanref{Frfp.Core.Kernel.pipeline\_compose}
\end{definition}

\begin{theorem}[Composition preserves typing]
If $\Pi_1 : X \to Y$ and $\Pi_2 : Y \to Z$, then
$(\Pi_2 \circ \Pi_1).\mathrm{source} = X$ and
$(\Pi_2 \circ \Pi_1).\mathrm{target} = Z$.
\leanref{Frfp.Core.Kernel.pipeline\_compose}
\end{theorem}

% -------------------------------------------------------------------
\subsection{Primitive Phase Partition}
% -------------------------------------------------------------------

\begin{definition}[Primitive phase]
The phase of a primitive is defined by:
\[
  \mathrm{phase}(p) =
  \begin{cases}
    \text{explicit}  & \text{if } p \in \{\RI, \EC, \ED\} \\
    \text{boundary}  & \text{if } p = \RB \\
    \text{tacit}     & \text{if } p \in \{\TE, \HFD\}
  \end{cases}
\]
\leanref{Frfp.Minimal.BasisSweep.phaseOf}
\end{definition}

\begin{theorem}[Phase partition]
The phase function is total: every primitive has exactly one phase.
\leanref{Frfp.Minimal.BasisSweep.phase\_partition}
\end{theorem}

\begin{theorem}[Boundary uniqueness by phase]
A primitive has boundary phase if and only if it is $\RB$.
\leanref{Frfp.Minimal.BasisSweep.phase\_boundary\_iff}
\end{theorem}

% -------------------------------------------------------------------
\subsection{The No-Return Principle}
% -------------------------------------------------------------------

\begin{theorem}[No tacit-to-explicit return]
There is no sequence of primitives that goes from $\Tobj$ to $\Eobj$:
\[
  \nexists\ p_1, \ldots, p_n : \mathrm{Primitive},\quad
  p_1.\mathrm{source} = \Tobj \wedge p_n.\mathrm{target} = \Eobj
\]
\end{theorem}

This follows from Theorem~1.1 (ETS) by induction: every primitive either stays
in tacit space ($\Tobj \to \Tobj$) or goes to $\Eobj$ from $\varnothing$ or
$\Eobj$.

The one-directional boundary
$\Eobj \xrightarrow{\RB} \Tobj$ is a genuine asymmetry in the categorical
structure: not a functor, not an equivalence, but a one-way transition that
formalizes the tacit grounding of AI outputs by human judgment.

% -------------------------------------------------------------------
\subsection{Immutability of the Kernel}
% -------------------------------------------------------------------

\begin{note}[Kernel immutability]
The definitions of \texttt{Object}, \texttt{Primitive},
\texttt{Primitive.source}, and \texttt{Primitive.target} are immutable. No module
in the library may shadow, redefine, or extend them.
\end{note}

This is enforced architecturally: \texttt{Kernel.lean} has no imports other than
Lean's standard library, and all other modules depend on it. Any change to
Kernel invalidates all 3305 downstream build jobs.
\leanref{Frfp.Core.Kernel --- no extends, no shadows}

\newpage
% ===================================================================
\section{The Reduced Axiom Basis}
\label{ch:basis}
% ===================================================================

\subsection{From 11 Premises to 7}

The original published FRFP theory enumerates 11 named premises: ETS, RB, AE, MD,
HEG, HEC, CP, IL, AR, NTER, CSC. The Lean formalization establishes that four of
these (ETS, RB, AE, MD) are theorems derivable directly from the kernel
definitions. The minimal independent basis therefore has exactly \textbf{7
primitive assumptions}.

\begin{center}
\begin{tabular}{llll}
\toprule
Premise & Original status & Lean status & Location \\
\midrule
ETS  & Axiom & \textbf{Theorem} (kernel typing) & Thm 2.1 \\
RB   & Axiom & \textbf{Theorem} (kernel enumeration) & Thm 2.2 \\
AE   & Axiom & \textbf{Theorem} (AI-executable typing) & Thm 2.3 \\
MD   & Axiom & \textbf{Theorem} (AI/human partition) & Thm 2.4 \\
HEG  & Axiom & \textbf{Axiom} (retained) & \S\ref{sec:core-irreducible} \\
HEC  & Axiom & \textbf{Axiom} (retained) & \S\ref{sec:core-irreducible} \\
CP   & Axiom & \textbf{Axiom} (retained) & \S\ref{sec:core-irreducible} \\
IL   & Axiom & \textbf{Axiom} (retained) & \S\ref{sec:interaction} \\
AR   & Axiom & \textbf{Axiom} (retained) & \S\ref{sec:interaction} \\
NTER & Axiom & \textbf{Axiom} (retained) & \S\ref{sec:interaction} \\
CSC  & Axiom & \textbf{Axiom} (retained) & \S\ref{sec:interaction} \\
\bottomrule
\end{tabular}
\end{center}

% -------------------------------------------------------------------
\subsection{Derivable Theorems (No Longer Axioms)}
% -------------------------------------------------------------------

\subsubsection*{ETS --- Explicit--Tacit Separation}

\begin{theorem}[ETS is a theorem]
There is no primitive from $\Tobj$ to $\Eobj$:
\[
  \forall p : \mathrm{Primitive},\
  \neg(p.\mathrm{source} = \Tobj \wedge p.\mathrm{target} = \Eobj)
\]
\leanref{Frfp.Minimal.ReducedBasis.ETS\_derivable} --- proved by delegating to
\texttt{no\_morphism\_T0\_to\_E0}.
\end{theorem}

\begin{note}
ETS holds because the primitive enumeration simply contains no $\Tobj \to \Eobj$
morphism. It is a fact about the closed enumeration, not an independent assumption.
\end{note}

\subsubsection*{RB --- Boundary Uniqueness}

\begin{theorem}[RB uniqueness is a theorem]
$\RB$ is the unique primitive from $\Eobj$ to $\Tobj$:
\[
  \forall p : \mathrm{Primitive},\
  (p.\mathrm{source} = \Eobj \wedge p.\mathrm{target} = \Tobj) \Rightarrow p = \RB
\]
\leanref{Frfp.Minimal.ReducedBasis.RB\_derivable}
\end{theorem}

\subsubsection*{AE --- AI-Explicit Restriction}

\begin{theorem}[AE is a theorem]
Every AI-executable primitive stays within explicit space:
\[
  \forall p,\
  \textit{is\_AI\_executable}(p)
  \Rightarrow (p.\mathrm{source} \in \{\varnothing, \Eobj\})
              \wedge p.\mathrm{target} = \Eobj
\]
\leanref{Frfp.Minimal.ReducedBasis.AE\_derivable}
\end{theorem}

\subsubsection*{MD --- Mode Discipline}

\begin{theorem}[MD is a theorem]
AI-executable primitives always target $\Eobj$; human-only primitives always
target $\Tobj$:
\[
  (\forall p,\ \textit{is\_AI\_executable}(p) \Rightarrow p.\mathrm{target} = \Eobj)
  \wedge
  (\forall p,\ \textit{is\_human\_only}(p) \Rightarrow p.\mathrm{target} = \Tobj)
\]
\leanref{Frfp.Minimal.ReducedBasis.MD\_derivable}
\end{theorem}

% -------------------------------------------------------------------
\subsection{Core Irreducible Premises}
\label{sec:core-irreducible}
% -------------------------------------------------------------------

The three core premises below are retained as genuine axioms. They capture
domain commitments that cannot be derived from the kernel structure alone.

\subsubsection*{HEG --- Human-Exclusive Grounding}

\begin{axiom_env}[Human-Exclusive Grounding (HEG)]
Tacit evaluation and human final decision --- the primitives $\TE$ and $\HFD$ ---
are exclusively executable by humans. No AI system can perform tacit grounding.

\emph{Formal skeleton}:
$\mathrm{HEG}_{\mathrm{skel}}:\ \TE.\mathrm{target} = \Tobj \wedge
\HFD.\mathrm{target} = \Tobj$

\leanref{Frfp.Minimal.ReducedBasis.HEG\_skeleton\_derivable}

\emph{Policy residue}: The normative claim --- that only humans may execute
$\TE$ and $\HFD$ --- is a governance premise.
\end{axiom_env}

\subsubsection*{HEC --- Human-Exclusive Closure}

\begin{axiom_env}[Human-Exclusive Closure (HEC)]
The final-decision primitive $\HFD$ is an endomorphism of tacit space that
only humans may apply.

\emph{Formal skeleton}:
$\mathrm{HEC}_{\mathrm{skel}}:\ \HFD.\mathrm{source} = \Tobj \wedge
\HFD.\mathrm{target} = \Tobj$

\leanref{Frfp.Minimal.ReducedBasis.HEC\_skeleton\_derivable}
\end{axiom_env}

\subsubsection*{CP --- Compositional Pipelines}

\begin{axiom_env}[Compositional Pipelines (CP)]
Pipelines compose in a way that preserves the source and target boundary:
\[
  (\Pi_2 \circ \Pi_1).\mathrm{source} = \Pi_1.\mathrm{source},\quad
  (\Pi_2 \circ \Pi_1).\mathrm{target} = \Pi_2.\mathrm{target}
\]
\leanref{Frfp.Minimal.ReducedBasis.CP\_skeleton\_derivable}
\end{axiom_env}

% -------------------------------------------------------------------
\subsection{Interaction Irreducible Premises}
\label{sec:interaction}
% -------------------------------------------------------------------

\subsubsection*{IL --- Intent Locking}

\begin{axiom_env}[Intent Locking (IL)]
All AI steps in a protocol trace share a common intent tag. An AI step cannot
unilaterally change the declared intent of the interaction:
\[
  \forall a, b \in \mathcal{P},\
  a.\mathrm{isAI} = \top \wedge b.\mathrm{isAI} = \top
  \Rightarrow a.\mathrm{intentTag} = b.\mathrm{intentTag}
\]
\emph{Source anchor}: Temporal safety / invariant-over-traces semantics.
\textbf{[Pnueli 1977]}
\end{axiom_env}

\subsubsection*{AR --- Ambiguity Resolution}

\begin{axiom_env}[Ambiguity Resolution (AR)]
Every ambiguous step in a protocol trace must be clarified before the
interaction proceeds:
\[
  \forall a \in \mathcal{P},\
  a.\mathrm{ambiguous} = \top \Rightarrow a.\mathrm{clarified} = \top
\]
\emph{Source anchor}: Liveness over traces; safety/liveness decomposition.
\textbf{[Alpern--Schneider 1985]}
\end{axiom_env}

\subsubsection*{NTER --- No Tacit Emulation/Reconstruction}

\begin{axiom_env}[No Tacit Emulation or Reconstruction (NTER)]
No AI step attempts to construct a morphism from $\Tobj$ to $\Eobj$, i.e.,
no AI step attempts to emulate or substitute for tacit human judgment.

\emph{Formal skeleton}:
$\forall p,\ p.\mathrm{source} = \Tobj \Rightarrow p.\mathrm{target} \neq \Eobj$

\leanref{Frfp.Minimal.ReducedBasis.NTER\_skeleton\_derivable}
\end{axiom_env}

\subsubsection*{CSC --- Correctness Stability Constraint}

\begin{axiom_env}[Correctness Stability Constraint (CSC)]
Once a step in a protocol trace is correctness-locked, all subsequent steps are
correctness-locked. Correctness locking is monotone:
\[
  (\exists a \in \mathcal{P},\ a.\mathrm{correctnessLocked} = \top)
  \Rightarrow
  \forall b \in \mathcal{P},\ b.\mathrm{correctnessLocked} = \top
\]
\emph{Source anchor}: Temporal safety (invariant preservation).
\textbf{[Pnueli 1977]}
\end{axiom_env}

% -------------------------------------------------------------------
\subsection{Source Anchoring of External Assumptions}
% -------------------------------------------------------------------

Modules beyond the kernel require additional axioms from established mathematics.

\begin{center}
\begin{tabular}{lll}
\toprule
Family & Representative axioms & Primary source \\
\midrule
IEEE 754 float arithmetic & Associativity, commutativity, order & [IEEE 754-2019] \\
Probability / measure & Countable additivity, normalization & [Billingsley 1995] \\
Stopping times & Measurability, Markov property & [Billingsley 1995] \S36 \\
Stochastic tail bounds & Markov inequality, geometric decay & [Durrett 2019] \\
Term rewriting / ARS & Termination, confluence & [Baader--Nipkow 1998] \\
Newman's Lemma & Confluence from termination & [Newman 1942] \\
Category / groupoid & Associativity, identity, inverses & [Mac Lane 1971] \\
Real analysis & Geometric limit-to-zero & [Rudin 1976] \\
Temporal safety & Safety properties over traces & [Pnueli 1977] \\
Safety vs liveness & Decomposition of liveness & [Alpern--Schneider 1985] \\
\bottomrule
\end{tabular}
\end{center}

\leanref{Frfp.Minimal.ReducedBasis.sourceBackedFamilies}

% -------------------------------------------------------------------
\subsection{Decomposition: Skeleton vs Policy Residue}
% -------------------------------------------------------------------

Each of the 7 retained axioms decomposes into two parts:
\begin{itemize}
  \item \textbf{Skeleton}: the mathematically structural part, derivable or
        source-anchored to published mathematics.
  \item \textbf{Policy residue}: the irreducible normative/governance commitment.
\end{itemize}

\begin{center}
\begin{tabular}{p{1cm}p{5.5cm}p{5.5cm}}
\toprule
Axiom & Skeleton & Policy residue \\
\midrule
HEG  & $\TE, \HFD$ target $\Tobj$ (Lean-proved) & Only humans may execute these \\
HEC  & $\HFD$ is a $\Tobj$-endomorphism (Lean-proved) & No AI substitution for closure \\
CP   & Composition preserves source/target (Lean-proved) & Pipelines are the only valid composition \\
IL   & Trace invariant (Pnueli 1977) & Intent is a shared, locked resource \\
AR   & Liveness (Alpern--Schneider 1985) & Ambiguity must be resolved \\
NTER & Structural = ETS (Lean-proved) & AI prohibited from tacit emulation attempts \\
CSC  & Safety/monotonicity (Pnueli 1977) & Correctness lock is irreversible \\
\bottomrule
\end{tabular}
\end{center}

The policy residue of all seven axioms is where the human-authority commitments
of FRFP live. These cannot be formalized as mathematical theorems --- they are
governance premises.

This is also the point at which human validation applies in the FRFP workflow.
Humans are not expected to re-verify Lean proof steps; they validate whether these
residues are acceptable commitments for the intended deployment context.

\leanref{Frfp.Minimal.ReducedBasis.DecompositionStatus}

\newpage
% ===================================================================
\section{Reduction and Confluence}
\label{ch:reduction}
% ===================================================================

\subsection{Overview}

This chapter defines the syntactic layer of FRFP --- the Task Decomposition Grammar
(TDG) that generates pipeline terms --- and the semantic reduction rules that
normalize them. We prove two structural properties: \emph{termination} (every
reduction sequence is finite) and \emph{confluence} (the order of reductions does
not affect the final result). Together these establish that every pipeline term
has a unique normal form modulo navigation equivalence.

% -------------------------------------------------------------------
\subsection{Task Decomposition Grammar}
% -------------------------------------------------------------------

\begin{definition}[TDG terms]
A TDG term over the kernel primitives is a well-typed finite sequence:
\[
  t ::= \varepsilon \mid p \cdot t
\]
where $p$ is a primitive and $t$ is a TDG term, subject to the typing
constraint that consecutive primitives compose.

TDG terms are syntactic objects. A TDG term \emph{denotes} a pipeline (morphism)
in $\Czero$ when interpreted by the semantic function $\llbracket \cdot \rrbracket$
(Chapter~\ref{ch:semantics}).
\end{definition}

\begin{definition}[Free category]
The free category $\mathbf{F}$ generated by the six primitives has:
\begin{itemize}
  \item Objects: $\{\varnothing, \Eobj, \Tobj\}$
  \item Morphisms: all well-typed finite compositions of primitives
\end{itemize}
$\mathbf{F}$ is the syntactic model; $\Czero$ is its semantic quotient.
\end{definition}

% -------------------------------------------------------------------
\subsection{Reduction Relations}
% -------------------------------------------------------------------

\begin{definition}[Explicit reduction]
A reduction relation $\to_e$ on TDG terms captures explicit-space
simplifications:
\[
  \EC \cdot \EC \to_e \EC,\qquad
  \ED \cdot \EC \to_e \ED
\]
together with similar rules for all combinations of
$\{\RI, \EC, \ED\}$.
\end{definition}

\begin{definition}[Tacit reduction]
A reduction relation $\to_t$ on TDG terms captures tacit-space simplifications:
\[
  \TE \cdot \TE \to_t \TE,\qquad
  \HFD \cdot \TE \to_t \HFD
\]
\end{definition}

\begin{definition}[Combined reduction]
The combined reduction $\to_c = \to_e \cup \to_t$, with contextual closure.
The reflexive-transitive closure is written $\twoheadrightarrow_c$.
\leanref{Frfp.Core.OperationalSemantics}, \leanref{Frfp.Core.SemanticCorrectness}
\end{definition}

% -------------------------------------------------------------------
\subsection{Normal Forms}
% -------------------------------------------------------------------

\begin{definition}[Normal form]
A TDG term $t$ is in \emph{normal form} if no reduction rule applies. A normal
form has the structure:
\[
  t_{\mathrm{nf}} =
  \underbrace{e_1 \cdots e_k}_{\text{explicit block}} \cdot
  \underbrace{\RB^?}_{\text{optional boundary}} \cdot
  \underbrace{h_1 \cdots h_m}_{\text{tacit block}}
\]
where each $e_i \in \{\RI, \EC, \ED\}$, $\RB^?$ is either $\RB$ or absent,
and each $h_j \in \{\TE, \HFD\}$.

\leanref{Frfp.Minimal.BasisSweep.PrimitiveTraceAdmissibleStrong}
\end{definition}

This three-block structure is the \textbf{protocol normal form}. It reflects the
operational semantics of an FRFP interaction: explicit work first, then
(optionally) handoff to human tacit evaluation.

% -------------------------------------------------------------------
\subsection{Termination}
% -------------------------------------------------------------------

\begin{theorem}[Termination]
Every reduction sequence in $\to_c$ is finite.
\leanref{Frfp.Core.ConfluenceProof --- termination component}
\end{theorem}
\begin{proof}
Define a measure $\mu(t)$ as the number of consecutive same-phase pairs in $t$.
Each reduction step strictly decreases $\mu(t)$. Since $\mu(t) \geq 0$, all
sequences terminate. \qed
\end{proof}

% -------------------------------------------------------------------
\subsection{Local Confluence}
% -------------------------------------------------------------------

\begin{definition}[Local confluence]
The relation $\to_c$ is \emph{locally confluent} if for any $t, t_1, t_2$ with
$t \to_c t_1$ and $t \to_c t_2$, there exists $t_3$ with
$t_1 \twoheadrightarrow_c t_3$ and $t_2 \twoheadrightarrow_c t_3$.
\end{definition}

\begin{theorem}[Local confluence of $\to_c$]
The combined reduction is locally confluent.
\leanref{Frfp.Core.ConfluenceProof.diamond\_property}
\end{theorem}
\begin{proof}
By case analysis on pairs of one-step reductions from $t$. Since explicit and
tacit reductions act on disjoint parts of the term (by the ETS theorem), any two
one-step reductions are either in the same phase at non-overlapping positions
(trivially confluent), in the same phase at overlapping positions (resolved by the
phase-specific confluence lemma), or in different phases (independently applicable
in either order). \qed
\end{proof}

\begin{note}
Lifted action commutativity --- navigation actions commute with explicit reduction
--- is axiomatized from the operational semantics.
\leanref{Frfp.Core.Confluence.lifted\_action\_commutes}

This property is implicit in the original published theory. The Lean formalization
makes it an explicit required lemma.
\end{note}

% -------------------------------------------------------------------
\subsection{Global Confluence (Church--Rosser)}
% -------------------------------------------------------------------

\begin{theorem}[Confluence]
By Newman's Lemma, termination together with local confluence implies global
confluence:
\[
  \forall t, t_1, t_2,\quad
  t \twoheadrightarrow_c t_1 \wedge t \twoheadrightarrow_c t_2
  \Rightarrow
  \exists t_3,\quad
  t_1 \twoheadrightarrow_c t_3 \wedge t_2 \twoheadrightarrow_c t_3
\]
\leanref{Frfp.Core.ConfluenceProof.confluence}\\
\emph{Source}: Newman's Lemma \textbf{[Newman 1942]}; abstract reduction systems
\textbf{[Baader--Nipkow 1998]} Ch.~2.
\end{theorem}

% -------------------------------------------------------------------
\subsection{Normal Form Uniqueness (Corrected)}
% -------------------------------------------------------------------

\begin{theorem}[Normal form uniqueness modulo navigation equivalence]
For any TDG term $t$, all normal forms reachable from $t$ are navigation-equivalent:
\[
  \forall t,\ \forall \mathrm{nf}_1, \mathrm{nf}_2,\quad
  t \twoheadrightarrow_c \mathrm{nf}_1
  \wedge t \twoheadrightarrow_c \mathrm{nf}_2
  \wedge \mathrm{IsNF}(\mathrm{nf}_1)
  \wedge \mathrm{IsNF}(\mathrm{nf}_2)
  \Rightarrow
  \mathrm{nf}_1 \navequiv \mathrm{nf}_2
\]
\leanref{Frfp.Core.SemanticCorrectness.unique\_nf\_mod\_nav}
\end{theorem}

\begin{correction}[Lean Correction]
The original published theory states normal forms are \emph{literally unique}
(equal). The Lean formalization shows they are unique only \emph{modulo navigation
equivalence} $\navequiv$. Two normal forms from the same term may differ in
navigation annotations while denoting the same pipeline.
\end{correction}

\begin{definition}[Navigation equivalence]
Two terms $t_1 \navequiv t_2$ if they differ only in navigation morphisms
(admissible path annotations added by the navigation layer). Navigation
morphisms do not change the computational content of a term.
\leanref{Frfp.Core.SemanticCorrectness.NavEquivalent}
\end{definition}

% -------------------------------------------------------------------
\subsection{Semantic Correctness}
% -------------------------------------------------------------------

\begin{theorem}[Observable invariance]
For any term $t$ and its normal form $\mathrm{nf}$, the observable correctness
judgment is preserved:
\[
  \obs(t) = \obs(\mathrm{nf})
\]
where $\obs = \gamma \circ \Sem \circ \mathrm{initial}$ is the composition of
semantic evaluation, correctness judgment extraction, and trajectory initialization.
\leanref{Frfp.Core.SemanticCorrectness.obsCorrect\_invariant}
\end{theorem}

\begin{definition}[Correctness judgment type]
The correctness judgment type has three values:
\[
  \mathrm{CorrectnessJudgment} =
  \{\mathrm{correct},\ \mathrm{incorrect},\ \mathrm{unknown}\}
\]
This is a new type made explicit by the Lean formalization; it was implicit in
the original paper.
\leanref{Frfp.Core.TacitDependence.CorrectnessJudgment}
\end{definition}

% -------------------------------------------------------------------
\subsection{Protocol Normal Form Theorem}
% -------------------------------------------------------------------

\begin{theorem}[Protocol normal form]
Any strongly admissible primitive trace decomposes as an explicit block,
optional boundary crossing, and tacit block:
\[
  \forall \mathcal{T},\
  \mathrm{AdmissibleStrong}(\mathcal{T})
  \Rightarrow
  \exists e_1\ldots e_k,\ h_1\ldots h_m,\ b \in \{\top, \bot\},
\]
\[
  \mathcal{T}
  = [e_1, \ldots, e_k]
    \mathbin{+\!\!+}
    (\text{if } b \text{ then } [\RB] \text{ else } [])
    \mathbin{+\!\!+}
    [h_1, \ldots, h_m]
\]
\leanref{Frfp.Minimal.BasisSweep.protocol\_normal\_form\_theorem} ---
Machine-checked proof.
\end{theorem}

\newpage
% ===================================================================
\section{Denotational and Operational Semantics}
\label{ch:semantics}
% ===================================================================

\subsection{Overview}

This chapter defines the semantic interpretation of FRFP terms. Two complementary
views are presented:

\begin{itemize}
  \item \textbf{Denotational semantics}: TDG terms denote morphisms in a
        Grothendieck construction over indexed categories.
  \item \textbf{Operational semantics}: pipeline execution is modeled as a
        relation on runs, connected to a stochastic trajectory space via the
        $\runToOmega$ bijection.
\end{itemize}

Both views contain significant structures made explicit by the Lean formalization
that were implicit or absent in the original published theory.

% -------------------------------------------------------------------
\subsection{The Grothendieck Construction}
% -------------------------------------------------------------------

FRFP uses a Grothendieck construction to model configurations --- the full state
of a pipeline execution including both categorical structure and index data.

\begin{definition}[Indexed category]
Define a contravariant functor:
\[
  \mathbf{E}_{\mathrm{idx}} : \mathbb{N}^{\mathrm{op}} \to \mathbf{Cat}
\]
mapping natural number indices $n$ to categories $\mathbf{E}_{\mathrm{idx}}(n)$
and index morphisms (decrements) to functors between them.
\end{definition}

\begin{definition}[Grothendieck object / configuration]
A configuration $\sigma$ is an object in the Grothendieck construction
$\int \mathbf{E}_{\mathrm{idx}}$:
\[
  \sigma = (n,\ X_n) \quad \text{where } n \in \mathbb{N},\
  X_n \in \mathrm{Ob}(\mathbf{E}_{\mathrm{idx}}(n))
\]
\leanref{Frfp.Core.Grothendieck.GrothendieckObject}
\end{definition}

\begin{definition}[Grothendieck morphism]
A morphism between configurations $(n, X_n) \to (m, X_m)$ consists of an index
morphism $n \to m$ and a morphism in the base category compatible with the
reindexing functor.
\leanref{Frfp.Core.Grothendieck.GrothendieckMorphism}
\end{definition}

% -------------------------------------------------------------------
\subsection{Lean Correction: Reduction Acts on Objects}
% -------------------------------------------------------------------

\begin{correction}[Lean Correction P-01]
The original published theory is ambiguous about whether reduction steps operate
on configurations, morphisms, or syntactic terms. The Lean formalization resolves
this: \textbf{reduction is a relation on configurations} (Grothendieck objects),
not on morphisms.
\end{correction}

\begin{definition}[Reduction step]
A reduction step is a relation
$\to_r \subseteq \mathrm{GrothendieckObject} \times \mathrm{GrothendieckObject}$.
\end{definition}

\begin{theorem}[Reduction induces morphism]
If $\sigma_1 \to_r \sigma_2$, then there exists a Grothendieck morphism $m$ with
$m.\mathrm{source} = \sigma_1$ and $m.\mathrm{target} = \sigma_2$:
\[
  \sigma_1 \to_r \sigma_2
  \Rightarrow
  \exists m : \mathrm{GrothendieckMorphism},\
  m.\mathrm{src} = \sigma_1 \wedge m.\mathrm{tgt} = \sigma_2
\]
\leanref{Frfp.Core.Grothendieck.reduction\_induces\_morphism}
\end{theorem}

% -------------------------------------------------------------------
\subsection{Semantic Functions}
% -------------------------------------------------------------------

\begin{definition}[Semantic evaluation]
The function $\Sem$ maps a configuration to an object in the kernel category:
\[
  \Sem : \mathrm{GrothendieckObject} \to \mathrm{Object}
\]
\leanref{Frfp.Core.TacitDependence.Sem}
\end{definition}

\begin{definition}[Correctness evaluation]
The function $\gamma$ maps a kernel object to a correctness judgment:
\[
  \gamma : \mathrm{Object} \to \mathrm{CorrectnessJudgment}
\]
\leanref{Frfp.Core.TacitDependence.gamma}
\end{definition}

\begin{definition}[Observable correctness]
The observable correctness function $\obs$ over runs:
\[
  \obs = \gamma \circ \Sem \circ \mathrm{initial}
       : \mathrm{Runs}(P) \to \mathrm{CorrectnessJudgment}
\]
\leanref{Frfp.Core.TacitDependence.obs}
\end{definition}

\begin{note}
$\Sem$, $\gamma$, $\obs$, and $\mathrm{CorrectnessJudgment}$ are all new
explicit types introduced by the Lean formalization. They were implicit in the
original paper's treatment of tacit-dependent correctness.
\end{note}

% -------------------------------------------------------------------
\subsection{Runs and Trajectories}
% -------------------------------------------------------------------

\begin{definition}[Run]
A run $\rho$ of program $P$ is a maximal execution trace. The type of runs of
$P$ is written $\mathrm{Runs}(P)$.
\end{definition}

\begin{definition}[Allowed stochastic trajectory]
The sample space $\Omega$ is the full space of stochastic trajectories. The
\emph{allowed} subspace $\Omega_{\mathrm{allowed}} \subseteq \Omega$ contains
exactly those trajectories consistent with a valid run of $P$.
\end{definition}

% -------------------------------------------------------------------
\subsection{The \texorpdfstring{$\runToOmega$}{runToOmega} Bijection}
% -------------------------------------------------------------------

\begin{correction}[Lean Correction P-02]
Remark A.104 in the original paper asserts a correspondence between runs and
allowed stochastic trajectories, but does not define it explicitly. The Lean
formalization requires a concrete injective map with a round-trip property.
\end{correction}

\begin{definition}[$\runToOmega$]
The run-to-trajectory correspondence:
\[
  \runToOmega : \{\rho \mid \mathrm{Runs}(P)\ \rho\} \to \Omega_{\mathrm{allowed}}
\]
with companion inverse:
\[
  \omegaToRun : \Omega_{\mathrm{allowed}} \to \{\rho \mid \mathrm{Runs}(P)\ \rho\}
\]
\leanref{Frfp.Core.OperationalSemantics.runToOmega}, \leanref{Frfp.Core.OperationalSemantics.omegaToRun}
\end{definition}

\begin{axiom_env}[$\runToOmega$ injectivity]
\[
  \runToOmega(\rho_1) = \runToOmega(\rho_2) \Rightarrow \rho_1 = \rho_2
\]
\leanref{Frfp.Core.OperationalSemantics.runToOmega_injective}
\end{axiom_env}

\begin{axiom_env}[$\runToOmega$ round-trip]
\[
  \omegaToRun(\runToOmega(\rho)) = \rho
\]
\leanref{Frfp.Core.OperationalSemantics.runToOmega_roundtrip}
\end{axiom_env}

\begin{note}
Injectivity and the round-trip property are axiomatized. A complete proof would
require formalizing the probability space construction. This is an open proof
obligation.

\textbf{Interpretation}: $\runToOmega$ is the formal bridge between the
operational semantics (discrete execution traces) and the probabilistic semantics
(measure on trajectory space). Without this bridge, probabilistic properties of
FRFP cannot be stated as theorems about runs.
\end{note}

% -------------------------------------------------------------------
\subsection{Denotational Semantics of Pipelines}
% -------------------------------------------------------------------

\begin{definition}[Pipeline denotation]
The denotational semantics $\llbracket \cdot \rrbracket$ maps a TDG term to its
morphism in $\Czero$:
\[
  \llbracket \varepsilon \rrbracket = \mathrm{id},\qquad
  \llbracket p \cdot t \rrbracket = \llbracket p \rrbracket \circ \llbracket t \rrbracket
\]
\end{definition}

\begin{theorem}[Denotational--operational agreement]
For any TDG term $t$ and its reduct $t'$, the denotations are equal:
\[
  t \to_c t' \Rightarrow \llbracket t \rrbracket = \llbracket t' \rrbracket
\]
\leanref{Frfp.Core.Semantics}
\end{theorem}

% -------------------------------------------------------------------
\subsection{Navigation Layer}
% -------------------------------------------------------------------

\begin{definition}[Navigation morphism]
A navigation morphism is an admissible path annotation on a TDG term, recording
which reachable states were visited during reduction. Navigation morphisms form a
groupoid.
\end{definition}

\begin{definition}[Navigation equivalence]
Two terms $t_1 \navequiv t_2$ if they are equal modulo navigation morphism
annotations.
\end{definition}

\begin{theorem}[Navigation equivalence is an equivalence relation]
\begin{itemize}
  \item Reflexivity: $t \navequiv t$.
  \item Symmetry: $t_1 \navequiv t_2 \Rightarrow t_2 \navequiv t_1$.
  \item Transitivity: $t_1 \navequiv t_2 \wedge t_2 \navequiv t_3
                      \Rightarrow t_1 \navequiv t_3$.
\end{itemize}
\leanref{Frfp.Core.SemanticCorrectness.nav\_equivalence\_is\_equivalence}
\end{theorem}

% -------------------------------------------------------------------
\subsection{Immutability of Pipelines}
% -------------------------------------------------------------------

\begin{theorem}[Pipeline immutability]
Executing a pipeline does not modify its structure. A pipeline's source, target,
and primitive are invariant under execution:
\[
  \mathrm{execute}(\Pi).\mathrm{structure} = \Pi.\mathrm{structure}
\]
\leanref{Frfp.Core.ImmutabilityProof}
\end{theorem}

This theorem captures a key FRFP invariant: pipelines are static specifications.
Execution produces outputs but does not rewrite the pipeline specification itself.

\newpage
% ===================================================================
\section{Dynamic and Probabilistic Layer}
\label{ch:dynamic}
% ===================================================================

\subsection{Overview}

This chapter covers two connected aspects of FRFP dynamics:

\begin{enumerate}
  \item \textbf{Epistemic state and credence degradation} --- how tacit knowledge
        degrades over time when not refreshed by human judgment.
  \item \textbf{Probabilistic stopping times and survival analysis} --- the
        stochastic model of pipeline resolution, connecting to the $\runToOmega$
        bijection.
\end{enumerate}

% -------------------------------------------------------------------
\subsection{Epistemic State}
% -------------------------------------------------------------------

\begin{definition}[Tacit state]
A tacit state records the current quality of human-grounded knowledge in a
pipeline execution. It is represented as a record with fields including a credence
value and a degradation parameter.
\end{definition}

\begin{definition}[Credence]
The credence of a tacit state is a rational number:
\[
  c \in [0, 1] \subset \mathbb{Q}
\]
\end{definition}

\begin{correction}[Lean Correction P-07]
The original paper uses credence notation without stating bounds explicitly.
The Lean formalization enforces bounds by construction:

\begin{lstlisting}
structure Credence where
  value  : Rat
  nonneg : 0 ≤ value
  le_one : value ≤ 1
\end{lstlisting}

\leanref{Frfp.Core.InstitutionalLayer.Credence}

Without these structural bounds, any arithmetic on credence could produce values
outside $[0,1]$, making the tacit degradation theorem unprovable.
\end{correction}

% -------------------------------------------------------------------
\subsection{Degradation}
% -------------------------------------------------------------------

\begin{definition}[Degradation operator $\delta$]
The degradation operator models the decay of tacit correctness over time when
not refreshed by human evaluation:
\[
  \delta : \mathrm{TacitState} \to \mathrm{TacitState}
\]
\end{definition}

\begin{definition}[Degradation formula --- corrected]
Given credence $c \in [0,1]$ and degradation rate $d \geq 0$, after $t$ steps
the credence is:
\[
  c'(t) = \max\!\Big(0,\ \min\!\Big(1,\ c - d \cdot t\Big)\Big)
\]
\end{definition}

\begin{correction}[Lean Correction P-07 (cont.)]
The original paper writes $c'(t) = c - d \cdot t$ without clamping. The Lean
formalization requires explicit clamping to keep credence within $[0,1]$; without
it, the non-negativity proof obligation fails.
\end{correction}

\begin{theorem}[Degradation monotonicity]
Credence is monotonically non-increasing under repeated degradation:
\[
  t_1 \leq t_2 \Rightarrow c'(t_2) \leq c'(t_1)
\]
\leanref{Frfp.Core.DynamicLayer.degradation\_monotone}
\end{theorem}

% -------------------------------------------------------------------
\subsection{Entropy Drift}
% -------------------------------------------------------------------

\begin{definition}[Entropy drift]
The entropy of the tacit component of a pipeline increases monotonically when no
human evaluation step occurs between AI steps:
\[
  \Delta H = H(t_2) - H(t_1) \geq 0 \quad
  \text{for } t_1 \leq t_2 \text{ with no HFD/TE between them}
\]
\end{definition}

\begin{theorem}[Entropy drift monotonicity]
Under continuous AI-only execution (no $\TE$ or $\HFD$ steps), the entropy of
the tacit component is non-decreasing.
\leanref{Frfp.Core.DynamicLayer.entropy\_drift\_monotone}
\end{theorem}

Entropy drift formalizes the FRFP claim that purely AI-driven pipelines without
human checkpoints progressively lose tacit validity. The $\TE$ primitive resets
the drift by incorporating human evaluation.

% -------------------------------------------------------------------
\subsection{Stopping Times}
% -------------------------------------------------------------------

\begin{definition}[Stopping time]
A stopping time $\tau$ for pipeline $P$ is a measurable random variable:
\[
  \tau : \Omega \to \mathbb{N} \cup \{\infty\}
\]
such that for all $n \in \mathbb{N}$, the event $\{\tau \leq n\}$ is measurable
with respect to the filtration generated by the first $n$ steps of the run.

\leanref{Frfp.Core.Probability.StoppingTime}
\end{definition}

\begin{definition}[Hazard rate]
The hazard rate $h(n)$ is the conditional probability that the pipeline resolves
at step $n$, given it has not resolved before step $n$:
\[
  h(n) = \mathbb{P}(\tau = n \mid \tau \geq n)
       = \frac{S(n-1) - S(n)}{S(n-1)}
\]
\end{definition}

\begin{definition}[Survival function]
The survival function $S(n)$ is the probability that the pipeline has not yet
resolved by step $n$:
\[
  S(n) = \mathbb{P}(\tau > n)
\]
\end{definition}

% -------------------------------------------------------------------
\subsection{Five Proven Probability Theorems}
% -------------------------------------------------------------------

The following five theorems were axiomatized in an earlier version of the
formalization (March 2026) and subsequently proved using Lean's Mathlib library
(April 2026). They are now machine-checked theorems, not assumptions.

\begin{theorem}[Survival product formula]
\[
  S(N) = \prod_{k=0}^{N} (1 - h(k))
\]
\leanref{Frfp.Core.ProbabilityProven.survival\_product\_formula}
\end{theorem}
\begin{proof}
By induction on $N$, using the definition of conditional probability and the
product rule for independent events. \qed
\end{proof}

\begin{theorem}[Hazard--survival relation]
\[
  h(n) = \frac{S(n-1) - S(n)}{S(n-1)}
\]
\leanref{Frfp.Core.ProbabilityProven.hazard\_survival\_relation}
\end{theorem}

\begin{theorem}[Survival monotonicity]
\[
  N \leq M \Rightarrow S(M) \leq S(N)
\]
\leanref{Frfp.Core.ProbabilityProven.survival\_monotone}
\end{theorem}

\begin{theorem}[Geometric survival]
When the hazard rate is constant, $h(n) = h$ for all $n$:
\[
  S(n) = (1-h)^{n+1}
\]
\leanref{Frfp.Core.ProbabilityProven.geometric\_survival}
\end{theorem}

\begin{theorem}[Infinite product limit]
When $h(k) > 0$ for all $k$:
\[
  \prod_{k=0}^{\infty} (1 - h(k)) = 0
\]
\leanref{Frfp.Core.ProbabilityProven.infinite\_product\_limit}
\end{theorem}

\begin{note}
The transition of these five theorems from axioms to machine-checked results
demonstrates that the FRFP probabilistic foundations are sound.
\end{note}

% -------------------------------------------------------------------
\subsection{Safe Horizon}
% -------------------------------------------------------------------

\begin{definition}[Safe horizon]
Given a degradation rate $d$ and an acceptability threshold $c_{\min}$, the safe
horizon $N^*(c, d)$ is the latest step at which the tacit state still has
acceptable credence:
\[
  N^*(c, d) = \max\{n \in \mathbb{N} \mid c'(n) \geq c_{\min}\}
            = \left\lfloor \frac{c - c_{\min}}{d} \right\rfloor
\]
\end{definition}

\begin{theorem}[Safe horizon characterization]
The safe horizon satisfies:
\[
  c'(N^*) \geq c_{\min} \wedge c'(N^*+1) < c_{\min}
\]
\leanref{Frfp.Core.DynamicLayer.safe\_horizon\_characterization}
\end{theorem}

The safe horizon determines when the pipeline \emph{must} invoke a $\TE$ or
$\HFD$ step --- i.e., when human evaluation is no longer optional.

% -------------------------------------------------------------------
\subsection{Tacit Irreversibility}
% -------------------------------------------------------------------

\begin{theorem}[Tacit irreversibility]
Once a pipeline transitions from explicit to tacit space (via $\RB$), this
transition cannot be reversed:
\[
  \nexists\ \text{sequence of primitives from } \Tobj \text{ to } \Eobj
\]
\leanref{Frfp.Core.TacitDependence}
\end{theorem}

This is a restatement of Theorem~1.7 (No-Return Principle) in dynamic terms:
once human evaluation is required, machine processing alone cannot restore
tacit validity.

\newpage
% ===================================================================
\section{Collective and Governance Layer}
\label{ch:collective}
% ===================================================================

\subsection{Overview}

This chapter extends FRFP from single-pipeline interactions to multi-agent
settings. Three layers build on the kernel:

\begin{enumerate}
  \item \textbf{Epistemic algebra} --- knowledge operators and belief revision
        for individual agents.
  \item \textbf{Institutional layer} --- population structures, credence
        distributions, and multi-agent impossibility results.
  \item \textbf{Governance layer} --- derived governance predicates and the
        safe extension theorem.
\end{enumerate}

% -------------------------------------------------------------------
\subsection{Epistemic Algebra}
% -------------------------------------------------------------------

\begin{definition}[Epistemic state]
An epistemic state for an agent is a record encoding the agent's current
knowledge about the pipeline's tacit validity. It includes:
\begin{itemize}
  \item A credence value $c \in [0,1]$.
  \item A belief set over possible pipeline configurations.
  \item An update history (finite sequence of knowledge operations).
\end{itemize}
\leanref{Frfp.Core.EpistemicAlgebra.EpistemicState}
\end{definition}

\begin{definition}[Knowledge operator]
A knowledge operator $K$ maps an epistemic state to a proposition:
\[
  K : \mathrm{EpistemicState} \to \mathrm{Prop}
\]
\leanref{Frfp.Core.EpistemicAlgebra.KnowledgeOp}
\end{definition}

\begin{definition}[Belief revision]
A belief revision operation updates an epistemic state given new evidence:
\[
  \mathrm{BeliefRevision} : \mathrm{EpistemicState} \times \mathrm{Evidence}
                           \to \mathrm{EpistemicState}
\]
\end{definition}

\begin{axiom_env}[Belief revision idempotence]
Revising with the same evidence twice is the same as revising once:
\[
  \mathrm{BeliefRevision}(\mathrm{BeliefRevision}(s, e), e)
  = \mathrm{BeliefRevision}(s, e)
\]
\leanref{Frfp.Core.EpistemicAlgebra.belief\_revision\_idempotent}
\end{axiom_env}

% -------------------------------------------------------------------
\subsection{Institutional Layer}
% -------------------------------------------------------------------

\begin{definition}[Population]
A population is a finite set of agents $\mathcal{A} = \{a_1, \ldots, a_n\}$
each carrying an epistemic state.
\leanref{Frfp.Core.InstitutionalLayer}
\end{definition}

\begin{definition}[Population credence]
The population credence $\bar{c}$ is the weighted average of individual
credences:
\[
  \bar{c} = \frac{1}{|\mathcal{A}|}\sum_{a \in \mathcal{A}} c(a)
\]
\end{definition}

\begin{theorem}[Population collapse bound]
Under continuous degradation without human evaluation, population credence
reaches zero in finite time:
\[
  \exists N,\ \bar{c}'(N) = 0
\]
\leanref{Frfp.Core.CollectiveLayer.population\_collapse\_bound}
\end{theorem}

% -------------------------------------------------------------------
\subsection{Consensus Impossibility}
% -------------------------------------------------------------------

\begin{theorem}[No grounding by consensus]
A consensus operation over AI agents cannot produce tacit grounding. Formally,
no combination of AI-only operations yields a correctness judgment:
\[
  \mathrm{ConsensusAI}(\mathcal{A})
  \not\Rightarrow
  \mathrm{CorrectnessJudgment.correct}
\]
\leanref{Frfp.Core.CollectiveLayer.consensus\_no\_grounding}
\end{theorem}

This theorem formalizes a key FRFP claim: voting, averaging, or ensemble
aggregation among AI systems cannot substitute for human tacit judgment.

\begin{note}
The theorem says explicit-space consensus (no matter how sophisticated) does not
produce a \texttt{correct} CorrectnessJudgment. The latter requires a $\HFD$
primitive, which is human-exclusive. The proof follows from the ETS theorem: no
AI operation produces a morphism targeting $\Tobj$.
\end{note}

% -------------------------------------------------------------------
\subsection{Governance Framework}
% -------------------------------------------------------------------

\begin{definition}[Strong IL --- global intent consistency]
\[
  \mathrm{IL}_{\mathrm{global}}(\mathcal{T})
  \iff
  \forall a, b \in \mathcal{T},\
  a.\mathrm{isAI} = \top \wedge b.\mathrm{isAI} = \top
  \Rightarrow
  a.\mathrm{intentTag} = b.\mathrm{intentTag}
\]
\leanref{Frfp.Minimal.BasisSweep.IL\_global}
\end{definition}

\begin{definition}[Strong AR --- full ambiguity resolution]
\[
  \mathrm{AR}_{\mathrm{global}}(\mathcal{T})
  \iff
  \forall a \in \mathcal{T},\
  a.\mathrm{ambiguous} = \top \Rightarrow a.\mathrm{clarified} = \top
\]
\leanref{Frfp.Minimal.BasisSweep.AR\_global}
\end{definition}

\begin{definition}[Strong CSC --- correctness lock propagation]
\[
  \mathrm{CSC}_{\mathrm{global}}(\mathcal{T})
  \iff
  \Big(\exists a \in \mathcal{T},\ a.\mathrm{locked} = \top\Big)
  \Rightarrow
  \forall b \in \mathcal{T},\ b.\mathrm{locked} = \top
\]
\leanref{Frfp.Minimal.BasisSweep.CSC\_global}
\end{definition}

\begin{definition}[Strong governance admissibility]
\[
  \mathrm{GA}(\mathcal{T})
  \iff
  \mathrm{IL}_{\mathrm{global}}(\mathcal{T})
  \wedge
  \mathrm{AR}_{\mathrm{global}}(\mathcal{T})
  \wedge
  \mathrm{CSC}_{\mathrm{global}}(\mathcal{T})
\]
\leanref{Frfp.Minimal.BasisSweep.GovernanceAdmissibleStrong}
\end{definition}

% -------------------------------------------------------------------
\subsection{Step Compatibility}
% -------------------------------------------------------------------

\begin{definition}[Step compatibility]
A new step $s$ is compatible with a strongly admissible trace $\mathcal{T}$ if:
\begin{enumerate}
  \item \textbf{Intent compatibility}:
        $\forall a \in \mathcal{T},\ a.\mathrm{isAI} = \top \wedge
        s.\mathrm{isAI} = \top \Rightarrow
        a.\mathrm{intentTag} = s.\mathrm{intentTag}$
  \item \textbf{Ambiguity compatibility}:
        $s.\mathrm{ambiguous} = \top \Rightarrow s.\mathrm{clarified} = \top$
  \item \textbf{Lock forward}:
        $(\exists a \in \mathcal{T},\ a.\mathrm{locked} = \top)
        \Rightarrow s.\mathrm{locked} = \top$
  \item \textbf{Lock backward}:
        $s.\mathrm{locked} = \top \Rightarrow
        \forall a \in \mathcal{T},\ a.\mathrm{locked} = \top$
\end{enumerate}
\leanref{Frfp.Minimal.BasisSweep.StepCompatibleStrong}
\end{definition}

% -------------------------------------------------------------------
\subsection{Governance Safe Extension Theorem}
% -------------------------------------------------------------------

\begin{theorem}[Governance safe extension]
Strong governance admissibility is preserved under compatible one-step extension:
\[
  \mathrm{GA}(\mathcal{T}) \wedge \mathrm{Compatible}(\mathcal{T}, s)
  \Rightarrow
  \mathrm{GA}(\mathcal{T} \mathbin{+\!\!+} [s])
\]
\leanref{Frfp.Minimal.BasisSweep.governance\_safe\_extension\_theorem}
--- machine-checked proof.
\end{theorem}
\begin{proof}
By case analysis on the membership of an arbitrary element of
$\mathcal{T} \mathbin{+\!\!+} [s]$: it is either in $\mathcal{T}$ (where
$\mathrm{GA}(\mathcal{T})$ applies) or equals $s$ (where the compatibility
conditions apply). \qed
\end{proof}

This theorem is the key safety result for FRFP governance: governance compliance
can be checked incrementally without re-auditing the entire trace history.

% -------------------------------------------------------------------
\subsection{Entropy Drift in Collective Settings}
% -------------------------------------------------------------------

\begin{theorem}[Collective entropy drift]
Under AI-only rounds (no $\TE$ or $\HFD$ steps from any agent), the
population-level entropy is non-decreasing:
\[
  H_{\mathrm{pop}}(t_2) \geq H_{\mathrm{pop}}(t_1) \quad
  \text{for } t_1 \leq t_2 \text{ with no human steps}
\]
\leanref{Frfp.Core.CollectiveLayer}
\end{theorem}

% -------------------------------------------------------------------
\subsection{Policy Residue in the Collective Layer}
% -------------------------------------------------------------------

\begin{center}
\begin{tabular}{p{1.8cm}p{5cm}p{5cm}}
\toprule
Property & Mathematical skeleton & Policy residue \\
\midrule
IL & Intent tag invariant over traces & Intent is a shared governance resource \\
AR & Every ambiguity flag implies clarification flag & Ambiguity is not acceptable in production \\
CSC & Correctness lock is trace-monotone & Once locked, correctness is never revisited \\
HEG/HEC & Tacit primitives target $\Tobj$ & Only humans may ground these \\
Consensus impossibility & No AI primitive targets $\Tobj$ & AI consensus cannot substitute for human judgment \\
\bottomrule
\end{tabular}
\end{center}

\newpage
% ===================================================================
\section{Novel Results from Formalization}
\label{ch:results}
% ===================================================================

\subsection{Overview}

Three non-trivial theorems emerge directly from the Lean formalization --- results
not present in the original published theory and arising solely from the proof
obligations imposed by the formal system. This chapter presents these theorems,
their proofs, their interpretations, and the open questions they suggest.

Additionally, 8 direct consequence theorems from the reduced basis are catalogued.

In FRFP terms, these results are machine-verified propositions whose publication
interpretation is still human-governed. That is: Lean certifies derivability;
humans decide commitment scope.

% -------------------------------------------------------------------
\subsection{Novel Theorem 1: First Unresolved Witness}
% -------------------------------------------------------------------

\begin{theorem}[First unresolved witness]
If an unresolved ambiguity exists anywhere in a protocol trace, there is a
\emph{canonical first} such occurrence:
\begin{align*}
  &\forall \mathcal{T},\
  \Big(\exists n,\
  \mathcal{T}[n].\mathrm{ambiguous} = \top
  \wedge \mathcal{T}[n].\mathrm{clarified} = \bot\Big)
  \Rightarrow \\
  &\exists n^*,\
  \Big(\mathcal{T}[n^*].\mathrm{ambiguous} = \top
  \wedge \mathcal{T}[n^*].\mathrm{clarified} = \bot\Big)
  \wedge \\
  &\quad \Big(\forall m < n^*,\
  \neg(\mathcal{T}[m].\mathrm{ambiguous} = \top
      \wedge \mathcal{T}[m].\mathrm{clarified} = \bot)\Big)
\end{align*}
\leanref{Frfp.Minimal.BasisSweep.first\_unresolved\_witness}
--- machine-checked proof.
\end{theorem}
\begin{proof}
Let $P(n) = \text{"step } n \text{ is unresolved"}$. By hypothesis, $\exists n,
P(n)$. By Lean's \texttt{Nat.find}, there exists a minimal witness $n^* =
\mathrm{Nat.find}(P)$ satisfying $P(n^*)$, and for all $m < n^*$, $\neg P(m)$.
\qed
\end{proof}

\paragraph{Why novel.} The original AR axiom says "all ambiguities are resolved"
--- a global property. This theorem localizes AR violations: when AR fails, it
fails at a specific first step. This is essential for:
\begin{itemize}
  \item \textbf{Audit trails}: governance audits can find the exact first violation
        point.
  \item \textbf{Incremental checking}: a trace is AR-compliant up to step $n$ iff
        $n < n^*$.
  \item \textbf{Repair protocols}: a minimal repair targets only step $n^*$, not
        the entire trace.
\end{itemize}

% -------------------------------------------------------------------
\subsection{Novel Theorem 2: Governance Safe Extension}
% -------------------------------------------------------------------

\begin{theorem}[Governance safe extension]
Strong governance admissibility (IL, AR, CSC) is preserved under compatible
one-step trace extension:
\[
  \mathrm{GA}(\mathcal{T}) \wedge \mathrm{Compatible}(\mathcal{T}, s)
  \Rightarrow
  \mathrm{GA}(\mathcal{T} \mathbin{+\!\!+} [s])
\]
\leanref{Frfp.Minimal.BasisSweep.governance\_safe\_extension\_theorem}
--- machine-checked proof.
\end{theorem}
\begin{proof}
Each of the three conjuncts of $\mathrm{GA}$ is preserved independently:
\begin{description}
  \item[\textbf{IL}:] For any two AI steps $a, b$ in $\mathcal{T}
        \mathbin{+\!\!+} [s]$, if both are in $\mathcal{T}$ then
        $\mathrm{IL}_{\mathrm{global}}(\mathcal{T})$ applies; if one is $s$ and
        the other is in $\mathcal{T}$, the intent compatibility condition applies.
  \item[\textbf{AR}:] For any ambiguous step in $\mathcal{T}
        \mathbin{+\!\!+} [s]$: if in $\mathcal{T}$,
        $\mathrm{AR}_{\mathrm{global}}(\mathcal{T})$ applies; if it is $s$, the
        ambiguity compatibility condition applies.
  \item[\textbf{CSC}:] The lock-forward and lock-backward conditions of
        compatibility ensure the lock propagates consistently. \qed
\end{description}
\end{proof}

\paragraph{Why novel.} The original FRFP governance axioms are stated as global
properties of a complete trace. Theorem 7.2 is an \emph{inductive} property ---
it says governance compliance is closed under one-step extension. This enables:
\begin{itemize}
  \item \textbf{Online governance checking}: evaluate each new step independently.
  \item \textbf{Compositional reasoning}: compliant sub-traces compose to a
        compliant whole.
  \item \textbf{Distributed protocols}: each agent can enforce local step
        compatibility; global compliance follows.
\end{itemize}

% -------------------------------------------------------------------
\subsection{Novel Theorem 3: Protocol Normal Form}
% -------------------------------------------------------------------

\begin{theorem}[Protocol normal form]
Every strongly admissible primitive trace decomposes into an explicit block, an
optional boundary crossing, and a tacit block:
\begin{align*}
  &\forall \mathcal{T} : \mathrm{List}(\mathrm{Primitive}),\
  \mathrm{AdmissibleStrong}(\mathcal{T}) \Rightarrow \\
  &\exists\ e_1\ldots e_k,\ h_1\ldots h_m,\ b \in \mathbb{B}, \\
  &\mathcal{T}
  = [e_1,\ldots,e_k]
    \mathbin{+\!\!+}
    (\text{if } b \text{ then } [\RB] \text{ else } [])
    \mathbin{+\!\!+}
    [h_1,\ldots,h_m]
\end{align*}
where each $e_i \in \Ecat = \{\RI, \EC, \ED\}$ and each
$h_j \in \Tcat \setminus \{\RB\} = \{\TE, \HFD\}$.

\leanref{Frfp.Minimal.BasisSweep.protocol\_normal\_form\_theorem}
--- machine-checked proof.
\end{theorem}

\paragraph{Why novel.} The three-block structure is mentioned informally in the
original theory. Theorem 7.3 makes it a formal theorem:
\begin{itemize}
  \item \textbf{Every} admissible trace has this structure --- it is not optional.
  \item The decomposition is canonical: the explicit block always precedes $\RB$,
        which always precedes the tacit block.
  \item The boundary $\RB$ appears \textbf{at most once} (by Theorem 1.2).
\end{itemize}

\paragraph{Implication for protocol design.} Any protocol that claims to be
FRFP-compliant but allows $\RB$ in a non-terminal position, or that interleaves
explicit and tacit steps, is not strongly admissible.

% -------------------------------------------------------------------
\subsection{Direct Consequences of the Reduced Basis}
% -------------------------------------------------------------------

\begin{theorem}[Tacit source classification]
$\forall p : \mathrm{Primitive},\ p.\mathrm{source} = \Tobj
\Rightarrow p \in \{\TE, \HFD\}$
\leanref{Frfp.Minimal.BasisConsequences.tacit\_source\_classification}
\end{theorem}

\begin{theorem}[Tacit source stays tacit]
$\forall p,\ p.\mathrm{source} = \Tobj \Rightarrow p.\mathrm{target} = \Tobj$
\leanref{Frfp.Minimal.BasisConsequences.tacit\_source\_stays\_tacit}
\end{theorem}

\begin{theorem}[AI steps never target tacit space]
$\forall p,\ \textit{is\_AI\_executable}(p) = \top \Rightarrow p.\mathrm{target} \neq \Tobj$
\leanref{Frfp.Minimal.BasisConsequences.ai\_never\_targets\_tacit}
\end{theorem}

\begin{theorem}[AI/human disjoint]
$\forall p,\ \neg(\textit{is\_AI\_executable}(p) \wedge \textit{is\_human\_only}(p))$
\leanref{Frfp.Minimal.BasisConsequences.ai\_human\_disjoint}
\end{theorem}

\begin{theorem}[Unique boundary entry]
$|\{p : \mathrm{Primitive} \mid p.\mathrm{source} = \Eobj \wedge p.\mathrm{target} = \Tobj\}| = 1$
\leanref{Frfp.Minimal.BasisConsequences.unique\_boundary\_entry}
\end{theorem}

\begin{theorem}[IL pairwise intent consistency]
Under IL, any two AI steps in the same trace share an intent tag.
\leanref{Frfp.Minimal.BasisConsequences.IL\_pairwise\_intent\_consistency}
\end{theorem}

\begin{theorem}[AR: no unresolved ambiguity]
Under AR, no ambiguous step remains unclarified.
\leanref{Frfp.Minimal.BasisConsequences.AR\_no\_unresolved\_ambiguity}
\end{theorem}

\begin{theorem}[CSC: lock is monotone]
Under CSC, the correctness-locked flag is monotone non-decreasing over the trace.
\leanref{Frfp.Minimal.BasisConsequences.CSC\_lock\_is\_monotone}
\end{theorem}

% -------------------------------------------------------------------
\subsection{Novelty Classification}
% -------------------------------------------------------------------

\begin{center}
\begin{tabular}{lll}
\toprule
Theorem & Type & How it arose \\
\midrule
7.1 First unresolved witness & New result & \texttt{Nat.find} forced minimal witness existence \\
7.2 Governance safe extension & New inductive property & Needed to prove $\mathrm{GA}$ is inductive \\
7.3 Protocol normal form & New completeness theorem & Admissibility definition crystallized three-block shape \\
Thms 7.4--7.11 & Direct consequences & Short proofs from basis definitions \\
\bottomrule
\end{tabular}
\end{center}

% -------------------------------------------------------------------
\subsection{Open Questions}
% -------------------------------------------------------------------

\begin{enumerate}[label=\textbf{Q\arabic*}]
  \item \textbf{Complete $\runToOmega$ proof.} The injectivity and round-trip
        properties are axiomatized. A full proof requires constructing a
        measurable bijection between the discrete run type and the
        measure-theoretic trajectory space.

  \item \textbf{Phase 2 primitives.} The current kernel is Phase 1 (6
        primitives). Can a Phase 2 extension be defined such that the analogues
        of ETS and RB uniqueness remain derivable theorems?

  \item \textbf{Mechanized minimality.} The minimality of the 7-axiom basis
        is established by manual witness construction. An automated model-finder
        could confirm or refute this mechanically.

  \item \textbf{Semantic grounding decidability.} The $\mathrm{CorrectnessJudgment}$
        type has value $\mathrm{unknown}$ for undecidable cases. Is there a
        formal characterization of when grounding is decidable vs.\ undecidable?

  \item \textbf{Probabilistic confluence.} Theorems 5.3--5.7 are deterministic
        properties of the hazard/survival functions. Is there a probabilistic
        confluence theorem: does the distribution over normal forms converge
        under repeated random reduction?
\end{enumerate}


% ===================================================================
\end{document}
